(1) 两侧均满足 \(\mu y_{tt}=Ty_{xx}\)，波数 \(k=\omega\sqrt{\mu/T}\)。点质量处位移连续 \(y(0^-,t)=y(0^+,t)=Y(t)\)，其横向受力由两侧张力分量差给出 \(M\ddot Y=T[y_x(0^+,t)-y_x(0^-,t)]\)。令左侧波为 \((e^{ikx}+r e^{-ikx})e^{-i\omega t}\)，右侧为 \(t e^{ikx}e^{-i\omega t}\)。
(2) 位移连续给 \(t=1+r\)。导数跳跃等于 \(ik[t-(1-r)]=2ikr\)，代入点质量动力学得 \(-M\omega^2t=2iTk r\)。联立得到 \(r=-M\omega^2/(M\omega^2+2iTk)\)、\(t=2iTk/(M\omega^2+2iTk)\)。即使质量是无耗散元件，反射和透射振幅也有非平凡复相位。
(3) 因两侧弦波阻抗相同，反射和透射能量率分别为 \(\mathcal R=|r|^2=M^2\omega^4/(M^2\omega^4+4T^2k^2)\)、\(\mathcal T=|t|^2=4T^2k^2/(M^2\omega^4+4T^2k^2)\)。直接相加得 \(\mathcal R+\mathcal T=1\)；质量可在周期中暂时储存动能，但周期平均没有净耗散。
(4) 因 \(k=\omega/c\)、\(c=\sqrt{T/\mu}\)，可把比值 \(M\omega^2/(Tk)\) 写成 \(M\omega c/T\)。当 \(\omega\to0\) 时该比值趋零，\(r\to0,t\to1\)，低频长波几乎不感受点质量。高频时惯性项 \(M\omega^2\) 占优，\(r\to-1,t\to0\)，质量趋于固定端。两个极限在波能量的比例上与不同物理边界相吻合。
